\begin{explanationbox}
  \textbf{\textcolor{CExplanation}{为什么会想到这个结论？}}

  遇到“对任意 $x$，某个二次式都 $>0$”，一个一个代入 $x$ 根本检查不完。
  但把二次式画成抛物线，只需看\textbf{最低处有没有落到 $x$ 轴上或下方}。
  因此，问题从“检查无数个数”变成了“检查一条曲线”。

  \medskip
  \textbf{\textcolor{CExplanation}{先用三个最简单的例子看懂}}

  \begin{itemize}[leftmargin=*,itemsep=0.3em]
    \item $x^2+1$：因为 $x^2\ge0$，所以 $x^2+1\ge1>0$，\textbf{恒正}。
    \item $x^2$：取 $x=0$ 得 $0$，所以\textbf{只能说恒非负}。
    \item $x^2-1$：取 $x=0$ 得 $-1$，所以\textbf{不恒正}。
  \end{itemize}
  这三个式子只有常数项不同，抛物线却分别\textbf{在轴上方、碰轴、穿轴}。

  \begin{center}
  \begin{tikzpicture}[x=1.18cm,y=0.72cm,>=stealth]
    \draw[->,gray] (-2,0)--(2.1,0) node[right] {$x$};
    \draw[->,gray] (0,-1.5)--(0,3.65) node[above] {$y$};
    \draw[green!60!black,very thick,domain=-1.45:1.45,samples=50] plot (\x,{\x*\x+1});
    \draw[orange!85!black,very thick,domain=-1.65:1.65,samples=50] plot (\x,{\x*\x});
    \draw[red!75!black,very thick,domain=-1.9:1.9,samples=50] plot (\x,{\x*\x-1});
    \fill[green!60!black] (0,1) circle(2pt);
    \fill[orange!85!black] (0,0) circle(2pt);
    \fill[red!75!black] (0,-1) circle(2pt);
    \node[anchor=west,green!45!black] at (2.25,2.6) {$y=x^2+1$（最低 $1$）};
    \node[anchor=west,orange!80!black] at (2.25,1.85) {$y=x^2$（最低 $0$）};
    \node[anchor=west,red!70!black] at (2.25,1.1) {$y=x^2-1$（最低 $-1$）};
  \end{tikzpicture}

  \end{center}

  \textbf{\textcolor{CExplanation}{为什么要检查两个条件？}}
  \begin{description}[style=nextline,leftmargin=2.8em,labelsep=0.5em,itemsep=0.45em]
    \item[\textbf{第一关：开口朝哪边？}]
      $a>0$ 时开口朝上，才可能整条在轴上方；$a<0$ 时开口朝下，往左右走得足够远，函数值一定会变成负数。
    \item[\textbf{第二关：有没有碰到横轴？}]
      令 $f(x)=0$，就是找抛物线和 $x$ 轴的交点：$\Delta>0$ 有两个交点，$\Delta=0$ 有一个交点，$\Delta<0$ 没有交点。要\textbf{始终大于零}，只接受最后一种。
  \end{description}

  \textbf{\textcolor{CExplanation}{把图像和计算连起来}}
  \[
    f(x)=a\left(x+\frac{b}{2a}\right)^2-\frac{\Delta}{4a}.
  \]
  当 $a>0$ 时，平方项不会小于 $0$，所以\textbf{最小值就是} $-\dfrac{\Delta}{4a}$。
  要让最低处还高于 $0$，必须 $-\dfrac{\Delta}{4a}>0$，这正好等价于 $\Delta<0$。

  \medskip
  \textbf{\textcolor{CExplanation}{考试什么时候用最快？}}
  题目出现“对任意实数 $x$，$ax^2+bx+c>0$”，尤其是\textbf{求参数范围}时，直接写
  \[
    \boxed{a>0\quad\text{且}\quad b^2-4ac<0},
  \]
  再解关于参数的不等式即可。\textbf{先确认二次项系数不为 $0$}；若它可能等于 $0$，要单独检查那种情况。
\end{explanationbox}
